\section{A residual for the two exterior truncations}

\begin{definition}[The two exterior truncations]
    \label{def:aklv_two_block_truncation}
    \lean{TNLean.MPS.AreaLaw.aklvRegionalHamiltonians,
        TNLean.MPS.AreaLaw.aklvTwoBlockHamiltonian}
    \leanok
    Let $0\le l\le r\le L$, let $k=(0,l,r,L)$, and let
    $0\le h_i\le\Id$ be supported on $\{i,i+1\}$.
    For $j\in\{0,1,2\}$ put
    \begin{align}
        A_j & =\sum_{\substack{k_j\le i                          \\i+1<k_{j+1}}}h_i, &
        J   & =\{i:i+1\in\{l,r\}\},     & Y & =\sum_{i\in J}h_i.
        \label{eq:aklv_two_block_regions}
    \end{align}
    Thus $\sum_i h_i=A_0+A_1+A_2+Y$ and $|J|\le2$.
    Choose real lower bounds $e_j\Id\le A_j$.
    For $T\subseteq\{0,1\}$ define $O_T=\{2j:j\in T\}$ and
    \begin{align}
        B_T     & =\sum_{j=0}^2(F_{j,T}-e_j\Id)+Y, &
        F_{j,T} & =
        \begin{cases}
            \min(A_j,(e_j+t)\Id), & j\in O_T,    \\
            A_j,                  & j\notin O_T.
        \end{cases}
        \label{eq:aklv_two_block_hamiltonian}
    \end{align}
    Here the minimum is defined by spectral calculus.
    Only the exterior sums are clipped; the middle sum and both connecting bonds
    are retained. This is the two-exterior geometry used in
    \cite[Section~6.1, after (5)]{Arad2013AreaLaw}.
\end{definition}

\begin{lemma}[Exponential residual on each partially truncated low-energy space]
    \label{thm:aklv_two_block_truncation_low}
    \lean{TNLean.MPS.AreaLaw.norm_aklv_two_block_truncation_low_le}
    \leanok
    \uses{def:aklv_two_block_truncation}
    Suppose $\varepsilon\ge0$ and $t\ge\varepsilon+5$.
    Each $B_T$ is Hermitian and
    $B_\varnothing=\sum_i h_i-(\sum_j e_j)\Id$.
    Let $G_T(\varepsilon)$ be the actual spectral subspace of $B_T$
    with eigenvalues at most $\varepsilon$.
    For every $T\subseteq\{0,1\}$,
    \begin{align}
        \bigl\|(B_\varnothing-B_{\{0,1\}})|_{G_T(\varepsilon)}\bigr\|
        \le64\exp\left(-\frac{t-\varepsilon}{8}\right).
        \label{eq:aklv_two_block_residual}
    \end{align}
\end{lemma}
\begin{proof}
    \leanok
    \uses{def:aklv_two_block_truncation,thm:consecutive_regional_truncation_low}
    The centered partial clippings and $Y$ are positive, so each $B_T$ is Hermitian.
    Partition the original bonds into the three regional interiors and $J$;
    this gives the untruncated identity. For each region, apply
    Lemma~\ref{thm:consecutive_regional_truncation_low} with
    $u=\varepsilon+5$ and imaginary time $1/8$.
    Writing $c=|J|\le2$, the resulting coefficient satisfies
    \begin{align*}
        \frac{e^{-(t-\varepsilon-5)/8-1}}{1/8}
        \frac{2c}{5-2c}
         & \le32e^{-3/8}e^{-(t-\varepsilon)/8}
        \le32e^{-(t-\varepsilon)/8}.
    \end{align*}
    The scalar shifts and the retained connecting terms cancel in
    $B_\varnothing-B_{\{0,1\}}$.
    Sum the two exterior residuals to obtain
    (\ref{eq:aklv_two_block_residual}).
\end{proof}

% This is an auxiliary residual on original nearest-neighbour data, not
% AKLV Theorem 6.1. Its source normalization, gap preservation, and normalized
% ground-vector comparison remain open; see the associated paper-gap note.
